% ============================================================
% 四川轻化工大学计算机科学与工程学院
% 本科毕业设计（论文）LaTeX 模板示例
% ============================================================

\documentclass{suse-cs-thesis}

% ---------- 封面信息 ----------
\thesistitle{事件驱动的分布式系统的研究}
\thesistitleen{Researches on the Scheme of Event-driven Distributed System}
\thesisauthor{张三}
\studentid{20220101001}
\major{软件工程}
\class{2022级1班}
\advisor{李四}
\thesisdate{二〇二六年五月}

% 设置校徽（PNG格式）
\logo{SUSE-icon.png}

% ---------- 正文开始 ----------
\begin{document}

% 封面
\makecover

% 独创性声明 + 授权说明（一页）
\makedeclaration

% 学术诚信与AI使用声明
\makeintegrity

% 中文摘要
\begin{cabstract}
本文研究了事件驱动的分布式系统的基本概念、分类、结构、特点和优缺点。分布式系统是一种软件系统，网络是其基础，但分布式系统的特点是其透明性。这种特性将分布式系统与普通网络使用区分开来。本文提出了一种事件驱动的分布式系统方案，该方案采用模块化方式实现，各模块相互协作以提高工作效率。利用这种模型可以快速实现分布式应用系统。最后，构建了一个简单的分布式模型来展示该工作框架并进行进一步讨论。

\ckeywords{分布式系统，事件驱动，网络，套接字}
\end{cabstract}

% 英文摘要
\begin{eabstract}
In this paper, the basic concept of the distributed system will be covered. It will also discuss the distributed system in foundation, classifications, structure, characteristics, advantages and shortage. The distributed system is a kind of software systems. Although network is its foundation, the distributed system is characterized by its transparency. This character distinguishes the distributed system from the ordinary use of the network. There are some discussions in this paper. In the following, a scheme of an event-driven distributed system will be issued. This model is issued by means of modularization and each block cooperates with each other in order to improve the work efficiency. A distributed application system can be quickly realized by utilizing such a kind of model. In the end, a simple distributed model will be constructed in order to show this working frame and further discussions.

\ekeywords{Distributed System, Event-driven, Network, Socket of Network}
\end{eabstract}

% 目录（确保所有目录页都无页眉）
\hideheader
\tableofcontents
\showheader
\clearpage

% 正文开始：固定行距22pt
\par\setlength{\baselineskip}{22pt}

% ---------- 第一章 绪论 ----------
\chapter{绪论}

文本文本文本文本文本文本文本文本文本文本文本文本文本文本文本文本文本文本文本文本文本文本文本。

\section{开发背景与意义}

文本文本文本文本文本文本文本文本文本文本文本文本文本文本文本文本文本文本文本文本文本文本文本。文本文本文本文本文本文本文本文本文本文本文本文本文本文本文本文本文本文本文本文本文本文本文本。

\section{标题}

文本文本文本文本文本文本文本文本文本文本文本文本文本文本文本文本文本文本文本文本文本文本文本。

\subsection{标题}

文本文本文本文本文本文本文本文本文本文本文本文本文本文本文本文本文本文本文本文本文本文本文本。

\subsection{标题}

文本文本文本文本文本文本文本文本文本文本文本文本文本文本文本文本文本文本文本文本文本文本文本。

\section{标题}

文本文本文本文本文本文本文本文本文本文本文本文本文本文本文本文本文本文本文本，用户表如表1-1所示。

% 表1-1 用户表数据表
\begin{table}[htbp]
    \centering
    \caption{用户表数据表}
    \label{tab:user}
    \begin{tabular}{lll}
        \toprule
        字段名称 & 数据类型 & 字段长度 \\
        \midrule
        ID & 自动编号 & 长整型 \\
        UserName & 文本 & 50 \\
        PWD & 文本 & 50 \\
        \bottomrule
    \end{tabular}
\end{table}

文本文本文本文本文本文本文本文本文本文本文本文本文本文本文本文本，Hotel数据表如表1-2所示。

% 表1-2 Hotel数据表
\begin{table}[htbp]
    \centering
    \caption{Hotel数据表}
    \label{tab:hotel}
    \begin{tabular}{lll}
        \toprule
        字段名称 & 数据类型 & 字段长度 \\
        \midrule
        ID & 自动编号 & 长整型 \\
        Name & 文本 & 50 \\
        Price & 数字 & 长整型 \\
        \bottomrule
    \end{tabular}
\end{table}

文本文本文本文本文本文本文本文本文本文本文本文本文本文本文本文本系统结构如图1-1所示。

% 图1-1 系统结构
\begin{figure}[htbp]
    \centering
    \fbox{\parbox{0.6\textwidth}{\centering\vspace{2cm}此处插入系统结构图\vspace{2cm}}}
    \caption{系统结构}
    \label{fig:structure}
\end{figure}

保证叙述中的指示与图的编号表示一致，图表分章节编号。图1-1表示第1章的第1个图，两数字中间用小段横线连接。图的标题在图下部，表的标题在表的顶部。图、表公式的编号各自独立，都从编号1开始。图表中的文字用五号字体。

论文中所有的表达式都按章节编号，每章的编号从第1个开始，"3-1"表示第3章第1个表达式。分章节编号的好处是其他章节编号改变不影响本章的编号顺序。

采用筛析法测定。用一套筛孔由大到小的标准筛，按顺序叠置在一起将被测定的树脂进行筛分，最后树脂主要落在一层筛面上，设该筛的筛孔尺寸为$d_1$，上一层筛子的筛孔为$d_2$，则按式（1-1）计算该层筛子上的树脂的尺寸\upcite{ref1}。

\begin{equation}
    d = \frac{d_1 + d_2}{2}
    \label{eq:resin}
\end{equation}

% ---------- 第二章 ----------
\chapter{标题}

相关内容没有添加，参照前面格式！

% ---------- 第三章 ----------
\chapter{标题}

相关内容没有添加，参照前面格式！

% ---------- 第四章 ----------
\chapter{标题}

相关内容没有添加，参照前面格式！

% ---------- 第五章 结论 ----------
\chapter{结论}

第一段落为结论部分，内容可根据摘要进行扩充。文本文本文本文本文本文本文本文本文本文本文本文本文本文本文本文本文本文本文本文本文本文本文本文本文本文本文本文本文本文本文本文本文本文本文本文本文本文本文本文本文本文本文本文本文本文本文本文本文本文本文本文本文本文本文本文本文本文本文本文本文本文本文本文本文本文本文本文本文本文本文本文本文本文本文本文本文本文本文本文本文本文本文本文本文本文本文本文本文本文本文本文本文本文本文本文本。

第二段落为展望部分，主要撰写本设计有哪些不足，这些不足下一步应该如何去完善。文本文本文本文本文本文本文本文本文本文本文本文本文本文本文本文本文本文本文本文本文本文本文本文本文本文本文本文本文本文本文本文本文本文本文本文本文本文本文本文本文本文本文本文本文本文本文本文本文本文本文本文本文本文本文本文本文本文本文本文本文本文本文本文本文本文本文本文本文本文本文本文本文本文本文本文本文本文本文本文本文本文本文本文本文本文本文本文本文本文本文本文本文本文本文本文本。

% ---------- 致谢 ----------
\begin{acknowledgements}

文本文本文本文本文本文本文本文本文本文本文本文本文本文本文本文本文本文本文本文本文本文本文本文本文本文本文本文本文本文本文本文本文本文本文本文本文本文本文本文本文本文本文本文本文本文本文本文本文本文本文本文本文本文本文本文本文本文本文本文本文本文本文本文本文本文本文本文本文本文本文本文本文本文本文本文本文本文本文本文本文本文本文本文本文本文本文本文本文本文本文本文本文本文本文本文本。

\end{acknowledgements}

% ---------- 参考文献 ----------
\begin{thebibliography}{99}

% 期刊
\bibitem{ref1} 陈家琪. 运动图像处理在车型识别中的应用[J]. 汽车工程, 2005, 20(6): 33-36.

% 期刊（3人以上）
\bibitem{ref2} 孙志军, 薛磊, 许阳明, 等. 深度学习研究综述[J]. 计算机应用研究, 2012, 29(8): 2806-2810.

% 期刊（3人）
\bibitem{ref3} 尹宝才, 王文通, 王立春. 深度学习研究综述[J]. 北京工业大学学报, 2015, (1): 48-59.

% 书籍
\bibitem{ref4} 陈家琪. C程序设计教程[M]. 北京: 新华出版社, 2004.

% 书籍（有版次）
\bibitem{ref5} 尼格尼维斯基. 人工智能:智能系统指南[M]. 第3版. 北京: 机械工业出版社, 2007.

% 学位论文
\bibitem{ref6} 李江. 红外图像人脸识别方法研究[D]. 长沙: 国防科技大学, 2005.

% 电子文献
\bibitem{ref7} 王明亮. 关于中国学术期刊标准化数据库系统工程的进展[EB/OL]. 1998-08-08[2018-12-11]. http://www.cajcd.edu.cn/pub/wml.txt/980810-2.html.

% 论文集
\bibitem{ref8} 汪学军. 中国农业转基因生物研发进展与安全管理[C]//中国国家生物安全框架实施国际合作项目研讨会论文集. 北京: 中国环境科学出版社, 2002: 22-25.

% 专利
\bibitem{ref9} 刘官山. 一种大数据学习系统: CN202210129842.1[P]. 2022-05-13.

\end{thebibliography}

% ---------- 附录 ----------
\appendix
\chapter{软件使用说明书}

\section{软件概述}

文本文本文本文本文本文本文本文本文本文本文本文本文本文本文本文本文本文本文本文本文本文本文本文本。

\subsection{功能}

文本文本文本文本文本文本文本文本文本文本文本文本文本文本文本文本文本文本文本文本文本文本文本文本。

\subsection{原理}

文本文本文本文本文本文本文本文本文本文本文本文本文本文本文本文本文本文本文本文本文本文本文本文本。

\section{软件安装}

文本文本文本文本文本文本文本文本文本文本文本文本文本文本文本文本文本文本文本文本文本文本文本文本。

\subsection{系统要求}

文本文本文本文本文本文本文本文本文本文本文本文本文本文本文本文本文本文本文本文本文本文本文本文本。

\subsection{安装前的准备}

文本文本文本文本文本文本文本文本文本文本文本文本文本文本文本文本文本文本文本文本文本文本文本文本。

\subsection{安装}

文本文本文本文本文本文本文本文本文本文本文本文本文本文本文本文本文本文本文本文本文本文本文本文本。

\subsection{安装后}

文本文本文本文本文本文本文本文本文本文本文本文本文本文本文本文本文本文本文本文本文本文本文本文本。

\section{运行说明}

文本文本文本文本文本文本文本文本文本文本文本文本文本文本文本文本文本文本文本文本文本文本文本文本。

\section{疑难解答}

文本文本文本文本文本文本文本文本文本文本文本文本文本文本文本文本文本文本文本文本文本文本文本文本。

\section{服务与支持}

文本文本文本文本文本文本文本文本文本文本文本文本文本文本文本文本文本文本文本文本文本文本文本文本。

\chapter{主要源程序}

\begin{verbatim}
//--------------------------------------------------------------------------
//挂上、卸下Keyboard hook，并运用自定义信息传递Hook 数据
//--------------------------------------------------------------------------
#include <vcl.h>
#pragma hdrstop
#include "Unit1.h"
//--------------------------------------------------------------------------
#pragma package(smart_init)
#pragma resource "*.dfm"
TKeyHookForm *KeyHookForm;
//--------------------------------------------------------------------------
__fastcall TKeyHookForm::TKeyHookForm(TComponent* Owner):TForm(Owner)
{
    inthook = 0;                           //计算按键次数变量归零
    FormStyle = fsStayOnTop;                //将KeyHookForm 维持在窗口最上面
    Button1->Enabled = true;                 //挂上和卸下Hook 的按钮状态
    Button2->Enabled = false;
}
//--------------------------------------------------------------------------
//挂上Hook 链，并设置定时器从共享内存中取回键盘信息
void __fastcall TKeyHookForm::Button1Click(TObject *Sender)
{
    inthook = 0;                           //计算按键次数变量归零
    //在keydll.dll 中，SetHook 函数执行注册Hook
    //行程到信息链中
    void (*SetHook)();                       //先声明有一SetHook 函数
    inst = LoadLibrary("keydll.dll");            //加载同一目录下的指定连接文件--keydll.dll;
    (FARPROC &)SetHook = ::GetProcAddress(inst,"SetHook");
    //取得dll 中SetHook 函数地址
    SetHook();                             //执行指向dll 中的SetHook 函数;
    Button1->Enabled = false;                 //挂上和卸下Hook的按钮状态
    Button2->Enabled = true;
}
//--------------------------------------------------------------------------
//卸下Hook 链，解除自定义的hook
void __fastcall TKeyHookForm::Button2Click(TObject *Sender)
{
    //在keydll.dll 中，RemoveHook 函数执行从信息链中卸下Hook 进程
    void (*RemoveHook)();                  //声明RemoveHook 函数
    inst = LoadLibrary("keydll.dll");           //加载同一目录下的指定连接文件--keydll.dll
    (FARPROC &)RemoveHook = GetProcAddress(inst,"RemoveHook");
    //取得dll 中RemoveHook 函数地址
    RemoveHook();                        //执行指向dll 中的RemoveHook 函数
    FreeLibrary(inst);                       //释放dll
    Button1->Enabled = true;                //挂上和卸下Hook 的按钮状态
    Button2->Enabled = false;
}
//--------------------------------------------------------------------------
//取得拦截自键盘中的自定义信息
void __fastcall TKeyHookForm::KeyHook(TMessage &Msg)
{
    //取得自定义信息中的按键名称
    char keytext[80];
    GetNameText(Msg.LParam, keytext, 80);
    AnsiString keystate;                     //取得自定义信息中的按键状态(检测用..意义不大...)
    keystate = GetKeyState((int)Msg.WParam);
    keystate = (keystate == "1")?"单击":"浮起";//C++的三元运算
    ListBox1->Items->Insert(0," 第 "+AnsiString(++inthook) + " 按键__" +
        AnsiString(keytext) + " ：状态>>" + keystate);         //由ListBox 组件显示拦截成果
}
//--------------------------------------------------------------------------
\end{verbatim}

\end{document}
